% !TeX root = ../CourseReport.tex
%=============================================================================
%  PART 4 --- GENAI SAFETY AND REFLECTION REPORT
%  PAGE BUDGET: 3 pages (excluding title, ToC)          Rubric weight: 10 pts
%
%  Required content:
%    - Safety of GenAI: safety considerations + risk-mitigation strategies
%    - Lessons learned and reflections (through the lens of your vision's
%      deliverables / KPIs / evaluations); what worked, what to improve.
%=============================================================================
\chapter{GenAI Safety and Reflection}
\label{ch:safety-reflection}

\section{Safety of GenAI}
% TODO: Safety considerations taken into account when using GenAI in Dishify and
% the strategies used to mitigate the risks. Good angles to cover:
%   - Dietary/allergy safety: restriction rules guardrails (restriction_rules.json)
%     that filter unsafe recipes before they reach the user.
%   - Prompt injection / untrusted input from voice and image ingestion.
%   - Hallucination in LLM reasoning and recipe augmentation, and how it is bounded.
%   - Privacy of user data (auth via Keycloak, what is sent to model providers).

\section{Lessons Learned and Reflections}
% TODO: Reflect through the lens of the key deliverables / KPIs / evaluations from
% the project vision. What worked well, what could be improved in future projects,
% and how this project shaped your understanding of GenAI and its applications.
